% Pillar-2 (ZVF) standalone introduction + related work
\section{Introduction}
\label{sec:p2-intro}
GRPO~\citep{guo2025deepseekr1, shao2024deepseekmath} learns from a group-relative
baseline: for each prompt it samples $G$ completions, scores them with a verifier,
and centers the rewards within the group. With binary correctness rewards, a group
whose completions are all correct or all incorrect has zero within-group variance
and zero centered advantage; under the submitted-runner assumptions, its centered
reward-contrast term is therefore zero. This does not assert that every
total-gradient contribution vanishes. As training proceeds
and the policy masters easy prompts or stalls on hard ones, the fraction of such
degenerate groups grows, and the effective batch that actually moves the policy
shrinks. We call this fraction the \emph{Zero-Variance Fraction} (ZVF) and ask how
well it diagnoses when GRPO stops learning.

ZVF is attractive because it is cheap, stack-agnostic, and mechanistically
motivated: it measures exactly the condition under which the group baseline nulls
that centered reward-contrast term. But a diagnostic is only as good as its predictive validity, and it is
easy to over-read a quantity that is mechanically entangled with reward sparsity,
group size, and accuracy. This paper therefore treats ZVF conservatively---as a
descriptive instrument whose reach we bound rather than a causal predictor.

We study ZVF as one pillar of \ours{}, a controlled benchmark of 70+ runs across
seven RL libraries and seven model families ($0.5$B--$1$T) on GSM8K
\citep{cobbe2021training}, HumanEval~\citep{chen2021codex}, and tool-use tasks.
Our contributions are: (i) measurements of ZVF across real runs---per-group
GSM8K reward tensors, a group-size sweep, scaling-law anchors, tool-use cells, and
a matched 8-seed comparison of GRPO with AERO, DAPO, Dr.~GRPO and GSPO in which
removing or reshaping zero-variance groups leaves held-out accuracy unchanged;
(ii) a signed decomposition showing why raw ZVF ranks healthy and unhealthy runs
wrongly; (iii) a gradient-utilization view, $\mathrm{GU} = 1 - \mathrm{ZVF}$,
linking the diagnostic to the fraction of the batch that carries signal; and (iv) a
cross-examination that identifies ZVF's central failure---aliasing mastery with
incapacity---and motivates magnitude- and sign-aware replacements.

\paragraph{Exact accounting, bounded interpretation.}
For a binary-reward prompt with success probability $p$, pass@$G$ is
$1-(1-p)^G$ and unanimous success is $p^G$. Their gap is therefore
$1-p^G-(1-p)^G=1-\mathrm{ZVF}$ under the stated iid-Bernoulli model. This is an
algebraic identity, not an empirical validation result. Three checked-in
per-group reward tensors independently recompute their reported ZVF values; no
separate 505-task identity-audit artifact is claimed. The identity tells us what
ZVF measures under the model---contrastive group yield---and what it does not:
held-out improvement, gradient direction, or causal value of an intervention.

\subsection{Related Work}
\label{sec:p2-related}
Reward sparsity and vanishing advantages are recurring obstacles in RL from
verifiable or human feedback~\citep{christiano2017deeprlhf, ouyang2022training, ahmadian2024backtobasics}. Variants of GRPO adjust the group baseline or advantage
normalization to preserve signal~\citep{lin2025cppo, gift2025, tong2025drgrpo}.
The diagnostic itself is not new: \citet{he2026advantage} define the Advantage
Collapse Rate (ACR)---the share of groups whose reward standard deviation falls
below a small threshold, i.e.\ ZVF---show that early ACR predicts final accuracy,
and use real-time ACR monitoring to drive AVSPO, which injects virtual reward
samples into homogeneous groups. \citet{nie2026gradstarvation} prove that the
degeneracy rate exceeds the i.i.d.\ Bernoulli prediction by Jensen's inequality.
Other methods act on zero-variance groups directly: by over-sampling and
filtering them (DAPO's dynamic sampling,~\citealp{yu2025dapo}; online difficulty
filtering,~\citealp{bae2025online}), by skipping prompts predicted to be
zero-variance (GRESO,~\citealp{zhang2026greso}), by adapting rollout allocation
(AERO,~\citealp{zhang2026aero}), by a variance-based curriculum
(VCRL,~\citealp{jiang2025vcrl}), or by extracting signal from them
(RL-ZVP,~\citealp{le2025rlzvp}).
Our aim is orthogonal---not to propose a new estimator but to bound what this
simple degeneracy statistic can and cannot tell a practitioner. In relating group reward variance
to the update, we connect to analyses that recast group-relative RL as an implicit
contrast or preference objective~\citep{wu2025grpodpo}, a link we make
explicit when discussing gradient utilization.
